\section{第四部分 总结}

\subsection{主要结论}
本作品完成了一套面向材料研发场景的双 RDK X5 嵌入式智能系统。它的复杂度来自跨域耦合：材料科学模型需要和 BPU/CPU 本地推理结合，AI 脑需要和 XRD/PL 表征线结合，具身脑需要和 LiDAR、深度、SLAM、Lab-FSD shadow 结合，机械臂工位和 F407 需要在安全边界内执行真实动作，公网平台需要把这些证据组织成评审可理解的页面。

相比只展示某个模型、某个网页或某台机械臂，本作品更接近一个实验室级嵌入式系统。它已经跑通的内容包括材料预测、四条表征线、SLAM 建图、Lab-FSD shadow、arm01 视觉门控取袋、F407 末端动作和公网证据平台；尚未完全闭环的内容也被明确标注为 shadow、降级或后续补齐。

\subsection{可扩展方向}
后续可扩展方向包括：
\begin{itemize}
  \item 补齐升降台限位和电磁铁驱动隔离，使 F407 工位具备更完整的无人值守安全保护；
  \item 优化 SLAM 地图质量，接入 Nav2 目标点和受限场景下的导航闭环；
  \item 扩展双臂十阶段剧本，将承碗、研磨、灌装、取烧结样和 XRD/PL 回填逐步真机化；
  \item 增加更多真实实验回流数据，改进 failure library、Fly-MB 记忆和材料预测置信度；
  \item 将公网展示站继续收敛为 Research Passport 和 Evidence Bundle，使复杂工程可复核、可答辩、可复现。
\end{itemize}

\subsection{心得体会}
项目推进中最重要的经验是：复杂嵌入式系统不能只追求“看起来全自动”，而要把每一层的真实边界写清楚。AI 预测不能替代烧制和表征，Lab-FSD 不能越过安全门直接接管底盘，F407 末端执行不能在限位未补齐时写成无人值守全闭环，公网平台不能开放危险控制。正是这些边界让系统在初赛阶段既能展示工程量，也能保持可信和可继续迭代。
